\chapter{Occupation Perturbations and the Weil Trace Comparison}
\label{v4-arith:w7-j:ATP-adelic-scattering}

Disjoint spectral intervals permit a full self-adjoint occupation
perturbation of the dilation generator. Its compact Fourier tests
have trace-class differences, even though the perturbation itself
is not trace class. The resulting spectral distribution is supported
on a half-line. This support excludes the full Weil zero functional.

\section{The scaling space and the occupation embedding}

\begin{definition}[idele norm coordinates]
\label{v4-arith:w7-j:def:H-idele}
Let $\mathbb A=\mathbb R\times\prod'_p\mathbb Q_p$ be the ring
of rational adeles, and put
$C_{\mathbb Q}=\mathbb Q^\times\backslash\mathbb A^\times$.
Choose the representative with positive real component and finite
components in $\widehat{\mathbb Z}^{\times}=\prod_p\mathbb Z_p^\times$.
The real component is the idele norm $r>0$. Use Haar measure
$dr/r\,d\mu_1$, with $\mu_1(\widehat{\mathbb Z}^{\times})=1$.
Put
\[
 U=L^2(\widehat{\mathbb Z}^{\times},\mu_1),\qquad e_0=1\in U,
 \qquad K=L^2(\mathbb R,d\tau)\otimes U.
\]
All tensor products of Hilbert spaces are completed.
The normalized Fourier transform in $t=\log r$ identifies the
scaling generator $-i\partial_t$ with $D_0=M_\tau\otimes1$.
Its domain is
\[
 \operatorname{Dom}D_0=
 \left\{v:\int_{\mathbb R}\tau^2\|v(\tau)\|_U^2d\tau<\infty\right\}.
\]
\end{definition}

To obtain the representative, multiply an idele by the reciprocal
of $\prod_pp^{v_p(a_p)}$, then by a sign. This makes every finite
valuation zero and the real component positive. Uniqueness follows
because a positive rational unit at every finite place is one.
The resulting product coordinates are continuous group coordinates.
Plancherel and the weak-derivative characterization of $H^1$
prove the generator and domain assertions on the full tensor product.
Functions constant on the second factor alone would not form a dense
domain in that tensor product.

Choose a real or complex $b\in C_c^\infty(0,1)$ with $\|b\|_2=1$.
For $m\ge1$, put $b_m(\tau)=b(\tau-m)$ and
$e_m=b_m\otimes e_0$. These are orthonormal. The map
$I:\ell^2(\mathbb N_{\ge1})\to K$, $Iv_m=e_m$, is an isometry.
For the occupation operator $Nv_m=(\log m)v_m$, use the maximal
domain $\sum_m(\log m)^2|a_m|^2<\infty$.
Its image gives
\[
 V_M=\sum_{m=1}^M(\log m)|e_m\rangle\langle e_m|,
 \qquad D_M=D_0+V_M.
\]
Inner products are linear in the first variable, so
$|e_m\rangle\langle e_m|v=\langle v,e_m\rangle e_m$.

\begin{proposition}[the full occupation perturbation]
\label{v4-arith:sc:full-domain}
The series
\[
 Vv=\sum_{m\ge1}(\log m)\langle v,e_m\rangle e_m
\]
defines a positive self-adjoint operator on the maximal domain
$\sum_m(\log m)^2|\langle v,e_m\rangle|^2<\infty$.
This domain contains $\operatorname{Dom}D_0$.
The operator $D=D_0+V$ is self-adjoint on $\operatorname{Dom}D_0$.
For every nonreal $z$, $(D_M-z)^{-1}$ converges strongly to
$(D-z)^{-1}$.
\end{proposition}
\begin{proof}
The diagonal series proves the assertion about $V$.
Decompose $K$ into the affected blocks
$K_m=L^2((m,m+1))\otimes\mathbb Ce_0$ and their common orthogonal
complement $K_{\mathrm{rest}}$. On $K_m$ write
\[
 A_m=M_\tau,\qquad
 B_m=A_m+(\log m)|e_m\rangle\langle e_m|.
\]
These are bounded self-adjoint operators with
$m\le A_m\le m+1$ and $m\le B_m\le m+1+\log m$.
The direct sum of $B_m$, together with $D_0$ on $K_{\mathrm{rest}}$,
is self-adjoint on its maximal direct-sum domain.
Since
\[
 m\|v_m\|\le\|A_mv_m\|\le(m+1)\|v_m\|,
 \qquad
 m\|v_m\|\le\|B_mv_m\|\le3m\|v_m\|,
\]
that domain equals $\operatorname{Dom}D_0$.
Bessel's inequality on each interval gives
$\|Vv\|^2\le\sum_m(\log m)^2\|v_m\|^2$.
Thus $V$ is defined there, and the direct sum equals $D_0+V$.

The resolvents of $D_M$ and $D$ agree on every fixed finite sum
of blocks once $M$ is sufficiently large. They agree on
$K_{\mathrm{rest}}$ for every $M$ and have norm at most $|\Im z|^{-1}$.
Finite block approximation of any vector proves strong convergence.
\end{proof}

The bounded perturbations satisfy
\[
 \|V_M\|=\log M,\quad
 \|V_M\|_1=\sum_{m\le M}\log m,\quad
 \|V_M\|_2^2=\sum_{m\le M}(\log m)^2.
\]
Thus $V$ is unbounded and is neither trace class nor Hilbert--Schmidt.
Its self-adjoint construction uses the specified interval embedding.
It does not identify that embedding with an arithmetic realization.

\section{Wave operators on the full model}

\begin{theorem}[existence of the occupation wave operators]
\label{v4-arith:w7-j:prop:wave-existence}
The strong limits
\[
 \Omega_\pm=\lim_{t\to\pm\infty}e^{itD}e^{-itD_0}
\]
exist on $K$ and are isometries. The same assertion holds with
$D_M$ in place of $D$.
\end{theorem}
\begin{proof}
Take $v$ in $C_c^\infty(\mathbb R)\otimes_{\mathrm{alg}}U$.
Only finitely many $e_m$ meet its spectral support. Each coefficient
of $Ve^{-itD_0}v$ is a constant multiple of
$\int\overline{b_m(\tau)}v_0(\tau)e^{-it\tau}d\tau$.
Two integrations by parts bound it by $C_v(1+|t|)^{-2}$.
Consequently $\int_{\mathbb R}\|Ve^{-itD_0}v\|dt<\infty$.
The common-domain identity
\[
 \frac d{dt}\bigl(e^{itD}e^{-itD_0}v\bigr)
 =ie^{itD}Ve^{-itD_0}v
\]
proves the Cauchy property at both ends of the real line.
The approximating operators are unitary. Their limits preserve
norms on this dense subspace and extend to isometries on all of $K$.
The finite sum gives the same proof for $D_M$.
\end{proof}

Completeness would additionally identify the ranges with the entire
absolutely continuous subspace of $D$. Existence alone does not make
that identification. If the two ranges are equal, then
$S=\Omega_+^*\Omega_-$ is unitary. Indeed both wave operators are
unitaries from $K$ onto their common closed range.

\section{A continuous trace map on Schwartz spectral tests}

\begin{lemma}[a block trace-norm estimate]
\label{v4-arith:sc:trace-bound}
For $h\in\mathscr S(\mathbb R)$ and every integer $N\ge1$,
\[
 \|h(B_m)-h(A_m)\|_1
 \le C_N(\log m)\sqrt{3+\log m}\,(1+m)^{-N}
 \max_{0\le j\le2}\sup_x(1+|x|)^N|h^{(j)}(x)|.
\]
Here the constant is independent of $m$ and $h$.
\end{lemma}
\begin{proof}
Choose a smooth cutoff $\chi_m$ equal to one on
$[m,m+1+\log m]$ and supported in $[m-1,m+2+\log m]$.
Fixed smooth transitions of width one give uniformly bounded
first and second derivatives. Put $q_m=\chi_m h$.
Functional calculus gives $q_m(A_m)=h(A_m)$ and
$q_m(B_m)=h(B_m)$.

Duhamel's formula and Fourier inversion give
\[
 \|q_m(B_m)-q_m(A_m)\|_1
 \le\frac{\log m}{2\pi}
       \int_{\mathbb R}|t|\,|\mathcal F_-q_m(t)|dt.
\]
For negative $t$, reverse the time integral in Duhamel's formula.
The same bound with $|t|$ follows. Cauchy--Schwarz with weight
$(1+t^2)^{-1}$ and Plancherel bound this integral by
$C(\|q_m'\|_2+\|q_m''\|_2)$.
The support length and uniform cutoff bounds give the displayed
estimate. For $m=1$ the difference is zero because $\log1=0$.
\end{proof}

\begin{theorem}[the occupation trace distribution]
\label{v4-arith:sc:trace-map}
The map
\[
 \mathcal T:\mathscr S(\mathbb R)\longrightarrow\mathcal S_1(K),
 \qquad \mathcal T(h)=h(D)-h(D_0)
\]
is continuous and linear, where $\mathcal S_1(K)$ has the trace norm.
Moreover,
\[
 T(h)=\operatorname{Tr}\mathcal T(h)
   =\sum_{m\ge1}\operatorname{Tr}(h(B_m)-h(A_m))
\]
defines a tempered distribution supported in $[1,\infty)$.
The trace differences for $D_M$ converge in trace norm to
$\mathcal T(h)$ for every Schwartz test.
\end{theorem}
\begin{proof}
Functional calculus respects the direct-sum decomposition.
On $K_{\mathrm{rest}}$ the difference is zero. On $K_m$ it is the
operator in Lemma~\ref{v4-arith:sc:trace-bound}.
For $N=3$ the sum of its bounds is finite, because
$\sum_m(\log m)\sqrt{3+\log m}(1+m)^{-3}<\infty$.
The sum therefore converges in trace norm and supplies a continuous
Schwartz seminorm bound. The finite truncations are its partial sums.
The trace is continuous in trace norm, which proves the formula
and temperedness. If $h$ is supported in $(-\infty,1)$, all affected
block functional calculi vanish. This proves the support assertion.
\end{proof}

\section{The support obstruction to an arithmetic trace identity}

\begin{theorem}[no lower-bounded tempered spectral realization]
\label{v4-arith:sc:support-obstruction}
There is no tempered distribution $T_0$ with support bounded below
such that
\[
 T_0(\mathcal F_+f)=W(f)\qquad(f\in C_c^\infty(\mathbb R)),
\]
where $W$ is the full xi zero functional.
In particular, the trace distribution of
Theorem~\ref{v4-arith:sc:trace-map} does not satisfy that identity.
\end{theorem}
\begin{proof}
Such an identity would extend $W$ to a tempered distribution in the
additive variable. On $C_c^\infty(0,\infty)$ the explicit formula gives
\[
 \nu_{1/2}(f)=-W(f)-\int_0^\infty
                \frac{e^{-t/2}}{e^{2t}-1}f(t)dt.
\]
Multiply by a smooth cutoff supported in $(0,\infty)$ and equal
to one for $t\ge1$. The density then decays exponentially.
The remainder of $\nu_{1/2}$ has compact support. Thus
$\nu_{1/2}$ would be tempered, and
Theorem~\ref{v4-arith:pm:centered} would imply RH.

Under RH, $W(f)=\mu_{\mathrm{ord}}(\mathcal F_+f)$.
Compact smooth functions are dense in the Schwartz space, and the
Fourier transform is its continuous automorphism. Equality on their
Fourier images therefore gives $T_0=\mu_{\mathrm{ord}}$ on all Schwartz tests.
There are infinitely many xi zeros, all in a fixed vertical strip.
Their ordinates are unbounded in absolute value. Conjugation gives
ordinates tending to negative infinity. The positive measure
$\mu_{\mathrm{ord}}$ therefore has support unbounded below, a contradiction.
\end{proof}

This argument excludes a proposed comparison for an explicitly
constructed operator and trace map. It does not prove or disprove RH.
An arithmetic realization must change the spectral carrier or the
trace construction and prove the resulting comparison.

\section{Scattering phases and the discrete spectral contribution}

\begin{proposition}[unitary scattering with a negative square trace]
\label{v4-arith:sc:unitary-counterexample}
On $K'=L^2(\mathbb R)\oplus\mathbb C$, let
$H_0=M_x\oplus0$ and $H_1=M_x\oplus1$ on the common maximal domain.
The wave operators are complete and their scattering operator is
the identity on the absolutely continuous subspace.
There is an entire Fourier square $h$ such that
\[
 \operatorname{Tr}(h(H_1)-h(H_0))<0.
\]
\end{proposition}
\begin{proof}
The absolutely continuous subspace is $L^2(\mathbb R)\oplus0$ for
both operators. On it $e^{itH_1}e^{-itH_0}$ is the identity for
every $t$, proving existence and completeness of both wave operators.
For a nonzero nonnegative $g\in C_c^\infty$, put
$\phi=\mathcal F_+g$ and $h=\phi\phi^\sharp$.
The functional-calculus difference is
$0\oplus(|\phi(1)|^2-|\phi(0)|^2)$.
Strict triangle inequality gives $|\phi(1)|<|\phi(0)|$, since
$e^{it}$ is not constant on any interval where $g>0$.
Thus the finite-rank trace is negative.
\end{proof}

For this pair the spectral-shift formula is the elementary identity
\[
 \operatorname{Tr}(h(H_1)-h(H_0))
 =\int_0^1h'(u)du.
\]
The real spectral-shift function is $\mathbf1_{(0,1)}$.
Its phase $\exp(-2\pi i\mathbf1_{(0,1)})$ is one almost everywhere.
The scattering phase loses this integer-valued discrete contribution.

\begin{proposition}[functional-equation phases]
\label{v4-arith:sc:phase}
The quotient $\xi(1/2+it)/\xi(1/2-it)$ equals one wherever defined
and has the constant continuation through its common zeros.
For an entire nonzero function $a$ satisfying
$a(\bar s)=\overline{a(s)}$, the quotient
$a(1/2+it)/a(1/2-it)$ has modulus one wherever defined.
\end{proposition}
\begin{proof}
The first assertion is the functional equation.
The second follows from conjugation. A common zero is isolated
unless the entire function vanishes identically. Consequently the
quotient, where unimodular, defines a unitary multiplication
operator on the real $L^2$ line after arbitrary values on its
measure-zero exceptional set. This does not identify it with wave
operators for a specified pair.
\end{proof}

Reflection symmetry $a(s)=\overline{a(1-\bar s)}$ alone has a
different meaning: it says that $a$ is real on the critical line.
For example, choose real even nonzero $c\in C_c^\infty(\mathbb R)$
and put
\[
 A(s)=\int c(t)e^{(s-1/2)t}dt,\qquad
 a(s)=(1+i(s-1/2))A(s).
\]
Then $a(s)=\overline{a(1-\bar s)}$, and it is entire of exponential
type with rapid decrease on closed vertical strips. Nevertheless
its quotient on the line is $(1-t)/(1+t)$ wherever $A$ is nonzero.
Reflection symmetry by itself does not give a unitary phase.

\section{The remaining arithmetic comparison}

\begin{definition}[the adele-class carrier]
\label{v4-arith:w7-j:def:connes-ac-space}
The adele-class quotient is $\mathbb Q^\times\backslash\mathbb A$,
where rational multiplication also acts on noninvertible adeles.
It differs from the idele-class group
$\mathbb Q^\times\backslash\mathbb A^\times$ used above.
A Hilbert or distributional realization on that quotient requires
its own space, action, domain, and trace functional.
\end{definition}

An arithmetic scattering comparison would specify self-adjoint
operators $A_0,A_1$ on a common Hilbert space and a trace map on
$\mathscr P$. If it uses an ordinary difference trace, it must prove
$h(A_1)-h(A_0)\in\mathcal S_1$ for every $h\in\mathscr P$ and
\[
 \operatorname{Tr}(h(A_1)-h(A_0))
 =P(f)-W_\infty(f)-W_{\mathrm{fin}}(f),\qquad h=\mathcal F_+f.
\]
If it uses a regularized or virtual trace, that functional and its
continuity must instead be constructed on the stated carrier.
The ordinary one-sided occupation model above fails the comparison
by a proved support obstruction. Scattering unitarity supplies
neither the missing trace identity nor positivity of its square class.
